\section{Introduction}

Nonlinear kernel SVMs expose a form of GPU workload that is easy to
misdiagnose. Their expensive RBF evaluations resemble dense linear algebra,
but training is not a sequence of independent matrix multiplications. A
decomposition solver repeatedly selects a working set, produces or reuses
kernel data, solves a box-constrained local quadratic program, updates a
global gradient, and uses the new state to make the next decision. The
performance of the system is therefore determined jointly by regular numeric
work and a state-dependent optimization trajectory.

This distinction becomes more important as the regular work is accelerated.
Tensor Cores make batches of dot products inexpensive, compact kernel values
increase cache coverage, and fused kernels reduce intermediate traffic. These
optimizations can expose a different limit: pairwise sequential minimal
optimization (SMO) commits one update before the next pair can be selected.
At that point, low occupancy is not itself the problem. It is evidence that
the algorithm has little independent work to schedule. Increasing occupancy
by speculating, batching, or replaying updates is useful only if the extra
work produces enough optimization progress.

We use RBF SVM training to study this interaction between representation,
data reuse, and dependency depth. Four observations organize the study.
First, precision is a property of the complete state transition, not of one
GEMM or cache in isolation: all producers and decision-making consumers of a
kernel value must observe a coherent representation. Second, global kernel
columns, local working-set rows, and dependent solver stages have different
costs and cannot be summarized by FLOP/s or occupancy. Third, the relevant
unit of efficiency is the optimization progress obtained from each expensive
piece of kernel data, rather than the number of pair updates executed. Fourth,
the profitable mechanism changes with feature dimension and training phase,
so a single static kernel is unlikely to be uniformly optimal.

These observations lead to three research questions.

\begin{itemize}
  \item \researchq{1}: How do kernel generation, reuse, and local optimization
        interact to determine end-to-end GPU SVM performance?
  \item \researchq{2}: Can a coherent low-precision kernel representation
        improve end-to-end time while retaining measured predictive quality?
  \item \researchq{3}: When does reducing the dependency depth of pair-SMO
        improve time to convergence, rather than merely moving work into
        additional kernel rows, columns, or dense solves?
\end{itemize}

The study makes four contributions. First, it introduces a cost decomposition
that separates full-data kernel-column misses, local-row consumption, and
dependent decision stages. This decomposition explains why superficially
similar ``kernel scans'' can differ by orders of magnitude. Second, it
develops coherent-FP16 execution, in which cached kernel values and the local
solver use the same rounded operator. On EMBER-65K, this design is
1.351$\times$ faster than a direct-FP32-cache control over seven balanced
fresh-process runs. Third, it characterizes the resulting bottleneck
migration: after kernel representation and reuse improve, the dependent local
solver becomes the largest measured GPU kernel. Fourth, it evaluates several
ways to shorten that chain and shows why reducing apparent depth is
insufficient when it increases kernel-data work or outer iterations.

The claims are intentionally bounded. The 1.351$\times$ result is a matched
comparison on one high-dimensional workload, with equal logical cache
capacity rather than equal cache bytes. A separate held-out evaluation finds
no statistically significant accuracy difference between the saved stock
FP32 and coherent-FP16 models in that experiment, but it does not establish
cross-workload equivalence. Coherent FP16 also does not satisfy the original
FP32-kernel KKT audit and is therefore a distinct low-precision optimization
contract, not an exact FP32 replacement. The solver experiments establish
boundaries for the measured designs; they do not preclude a parallel SVM
solver that discovers more productive optimization directions.
